\documentclass[12pt,a4paper]{ctexart}

\usepackage[a4paper,margin=2.6cm]{geometry}
\usepackage{amsmath,amssymb,amsthm,mathtools,bm}
\usepackage{physics}
\usepackage{esint}
\usepackage{booktabs}
\usepackage{enumitem}

\title{数学分析复习练习题}
\author{}
\date{}

\begin{document}

\maketitle

\begin{center}
\textbf{说明：}本练习题覆盖数学分析全部核心考点。
基础题约占 60\%，综合题约占 40\%。每题标注所考查的知识点。
\end{center}

\bigskip

% ==================== 第一章 引论 ====================
\section*{第一章\quad 引论}

\begin{enumerate}[leftmargin=*,label=\arabic*.]
  \item 设 $A = \{x \in \mathbb{R} \mid x^2 - 3x + 2 \le 0\}$，
    $B = \{x \in \mathbb{R} \mid 1 \le x \le 3\}$。
    求 $A \cup B$，$A \cap B$，$A \setminus B$。（考点：集合的运算）

  \item 判断函数 $f(x) = \ln(x + \sqrt{x^2+1})$ 的奇偶性。
    （考点：函数的奇偶性）

  \item 设 $f(x) = \frac{1}{1-x}$，求 $f(f(x))$ 和 $f(f(f(x)))$ 的表达式。
    （考点：复合函数）
\end{enumerate}

% ==================== 第二章 极限与连续 ====================
\section*{第二章\quad 极限与连续}

\begin{enumerate}[leftmargin=*,label=\arabic*.,start=4]
  \item 用 $\varepsilon$-$N$ 语言证明：$\displaystyle\lim_{n\to\infty}\frac{2n+1}{3n-2} = \frac{2}{3}$。
    （考点：数列极限的定义）

  \item 求极限 $\displaystyle\lim_{n\to\infty}\left(\frac{1}{\sqrt{n^2+1}} + \frac{1}{\sqrt{n^2+2}} + \cdots + \frac{1}{\sqrt{n^2+n}}\right)$。
    （考点：夹逼定理）

  \item 设 $a_1 = \sqrt{2}$，$a_{n+1} = \sqrt{2+a_n}$，证明 $\{a_n\}$ 收敛并求其极限。
    （考点：单调有界定理）

  \item 求极限 $\displaystyle\lim_{x\to 0}\frac{\tan x - \sin x}{x^3}$。
    （考点：等价无穷小与极限计算）

  \item 求极限 $\displaystyle\lim_{x\to 0}(1+3x)^{2/x}$。
    （考点：重要极限）

  \item 讨论函数 $f(x) = \frac{|x-1|}{x^2-1}$ 的间断点类型。
    （考点：间断点分类）

  \item 证明方程 $x^3 - 4x^2 + 1 = 0$ 在 $(0,1)$ 内至少有一个实根。
    （考点：零点存在定理）
\end{enumerate}

% ==================== 第三章 单变量函数的微分学 ====================
\section*{第三章\quad 单变量函数的微分学}

\begin{enumerate}[leftmargin=*,label=\arabic*.,start=11]
  \item 设 $f(x) = \begin{cases} x^2\sin\frac{1}{x}, & x \neq 0 \\ 0, & x = 0 \end{cases}$，
    讨论 $f$ 在 $x=0$ 处的可导性。（考点：导数定义）

  \item 求函数 $y = \ln\sin(x^2)$ 的导数。（考点：链式法则）

  \item 设 $y = x^2 e^x$，求 $y^{(n)}$。（考点：Leibniz 公式）

  \item 验证函数 $f(x) = x^3 - 3x$ 在 $[-1,1]$ 上满足 Rolle 定理的条件，
    并求出满足 $f'(\xi)=0$ 的 $\xi$。（考点：Rolle 定理）

  \item 利用 Lagrange 中值定理证明：当 $x > 0$ 时，
    $\frac{x}{1+x} < \ln(1+x) < x$。（考点：Lagrange 中值定理）

  \item 求极限 $\displaystyle\lim_{x\to 0}\frac{x-\sin x}{x^3}$。（考点：L'H\^opital 法则）

  \item 将 $f(x) = \ln(1+x)$ 在 $x=0$ 处展开至 $x^4$ 项（带 Peano 余项）。
    （考点：Taylor 公式）

  \item 求函数 $f(x) = x^3 - 6x^2 + 9x + 1$ 的极值点和拐点。
    （考点：导数的应用——极值与拐点）

  \item 求曲线 $y = \frac{x^2+1}{x-1}$ 的渐近线。
    （考点：函数的作图——渐近线）
\end{enumerate}

% ==================== 第四章 单变量函数的积分学 ====================
\section*{第四章\quad 单变量函数的积分学}

\begin{enumerate}[leftmargin=*,label=\arabic*.,start=20]
  \item 计算不定积分 $\displaystyle\int \frac{dx}{x^2-a^2}\ (a \neq 0)$。
    （考点：有理函数的积分）

  \item 计算不定积分 $\displaystyle\int e^x \sin x \, dx$。
    （考点：分部积分法）

  \item 计算定积分 $\displaystyle\int_0^{\pi/2} \sin^n x \, dx$（$n$ 为正整数）的递推公式。
    （考点：定积分的换元法）

  \item 利用换元法计算 $\displaystyle\int_0^1 \frac{\ln(1+x)}{1+x^2} dx$。
    （考点：定积分的换元）

  \item 计算广义积分 $\displaystyle\int_0^{+\infty} e^{-x^2} dx$。
    （考点：广义积分）

  \item 讨论广义积分 $\displaystyle\int_1^{+\infty} \frac{dx}{x^p}$ 的敛散性。
    （考点：广义积分的比较判别法）

  \item 求由曲线 $y = x^2$ 与 $y = \sqrt{x}$ 所围成图形的面积。
    （考点：定积分的应用——面积）

  \item 求曲线 $y = \frac{2}{3}x^{3/2}$ 在 $[0,3]$ 上的弧长。
    （考点：定积分的应用——弧长）
\end{enumerate}

% ==================== 第五章 常微分方程 ====================
\section*{第五章\quad 常微分方程}

\begin{enumerate}[leftmargin=*,label=\arabic*.,start=28]
  \item 求解微分方程 $\displaystyle\frac{dy}{dx} = \frac{y}{x} + \tan\frac{y}{x}$。
    （考点：齐次方程）

  \item 求解初值问题 $\begin{cases} y' + \frac{2}{x}y = \ln x, \\ y(1) = 0 \end{cases}$。
    （考点：一阶线性方程）

  \item 求方程 $y'' + 4y' + 3y = 0$ 的通解。
    （考点：常系数齐次线性方程）

  \item 求方程 $y'' + y = x\cos x$ 的一个特解。
    （考点：常系数非齐次线性方程——待定系数法）
\end{enumerate}

% ==================== 第六章 空间解析几何 ====================
\section*{第六章\quad 空间解析几何}

\begin{enumerate}[leftmargin=*,label=\arabic*.,start=32]
  \item 求过点 $P(1,2,-1)$ 且与平面 $2x - y + 3z + 1 = 0$ 平行的平面方程。
    （考点：平面方程）

  \item 求点 $M(2,1,5)$ 到直线 $\frac{x-1}{2} = \frac{y}{1} = \frac{z+3}{-1}$ 的距离。
    （考点：点到直线的距离）

  \item 指出曲面 $z = x^2 + y^2$ 的类型并描述其几何特征。
    （考点：二次曲面）
\end{enumerate}

% ==================== 第七章 多元函数微分学 ====================
\section*{第七章\quad 多元函数微分学}

\begin{enumerate}[leftmargin=*,label=\arabic*.,start=35]
  \item 讨论二重极限 $\displaystyle\lim_{(x,y)\to(0,0)}\frac{xy}{x^2+y^2}$ 是否存在。
    （考点：二重极限）

  \item 设 $z = e^{x^2+y^2}$，求 $\frac{\partial z}{\partial x}$，$\frac{\partial z}{\partial y}$，
    $\frac{\partial^2 z}{\partial x \partial y}$。
    （考点：偏导数）

  \item 设 $z = f(x^2-y^2, e^{xy})$，其中 $f$ 具有二阶连续偏导数，求
    $\frac{\partial z}{\partial x}$ 和 $\frac{\partial z}{\partial y}$。
    （考点：多元复合函数的链式法则）

  \item 求曲面 $z = x^2 + y^2$ 在点 $(1,1,2)$ 处的切平面和法线方程。
    （考点：曲面的切平面）

  \item 求函数 $f(x,y) = x^3 + y^3 - 3xy$ 的极值。
    （考点：多元函数的极值）

  \item 在 $x + y + z = 1$ 的条件下，求 $u = x^2 + y^2 + z^2$ 的最小值。
    （考点：条件极值——Lagrange 乘数法）
\end{enumerate}

% ==================== 第八章 多元函数积分学 ====================
\section*{第八章\quad 多元函数积分学}

\begin{enumerate}[leftmargin=*,label=\arabic*.,start=41]
  \item 计算二重积分 $\displaystyle\iint_D xy \,dxdy$，其中 $D$ 为由 $y=x$，$y=1$，$x=0$ 所围成的区域。
    （考点：二重积分的计算）

  \item 计算 $\displaystyle\iint_D e^{-(x^2+y^2)} dxdy$，其中 $D: x^2+y^2 \le R^2$。
    （考点：极坐标变换下的二重积分）

  \item 计算曲线积分 $\displaystyle\int_L (x^2-y^2)dx + 2xy\,dy$，其中 $L$ 为从 $(0,0)$ 到 $(1,1)$
    沿抛物线 $y=x^2$ 的弧段。（考点：第二类曲线积分的计算）

  \item 利用 Green 公式计算 $\displaystyle\oint_L (x^2y)dx - (xy^2)dy$，
    其中 $L$ 为圆周 $x^2+y^2 = a^2$ 的正向。（考点：Green 公式）

  \item 利用 Gauss 公式计算 $\displaystyle\oiint_S x\,dydz + y\,dzdx + z\,dxdy$，
    其中 $S$ 为球面 $x^2+y^2+z^2=a^2$ 的外侧。（考点：Gauss 公式）
\end{enumerate}

% ==================== 第九章 级数论 ====================
\section*{第九章\quad 级数论}

\begin{enumerate}[leftmargin=*,label=\arabic*.,start=46]
  \item 判断级数 $\displaystyle\sum_{n=1}^{\infty}\frac{n}{2^n}$ 的敛散性。
    （考点：正项级数——比值判别法）

  \item 判断级数 $\displaystyle\sum_{n=1}^{\infty}\frac{(-1)^{n-1}}{\sqrt{n}}$ 是绝对收敛、条件收敛还是发散。
    （考点：交错级数与绝对收敛）

  \item 求幂级数 $\displaystyle\sum_{n=1}^{\infty}\frac{x^n}{n}$ 的收敛域。
    （考点：幂级数的收敛半径）

  \item 将 $f(x) = x$（$-\pi < x < \pi$）展开为 Fourier 级数。
    （考点：Fourier 级数）
\end{enumerate}

\end{document}
